% ============================================================
%  Розширена секція оцінювання для mission-memory-paper
%  Замінює рядки 295--352 у main.tex
% ============================================================

% ============================================================
\section{Реалізація та експериментальна верифікація}
\label{sec:evaluation}
% ============================================================

\subsection{Платформа верифікації}
\label{sec:platform}

Архітектура верифікована на платформі юридичного штучного інтелекту LEX~AI~\citep{ovcharov2026edittrace}: 70+ MCP"=інструментів, 380M+ записів у pipeline даних, 1\,547 об'єднаних pull"=запитів за 105 днів роботи одного практика з LLM"=агентом.
Ця платформа є стрес"=тестом підсистеми пам'яті: один оператор з максимальною пропускною здатністю генерує щільний потік рішень, що моделює навантаження автономної місії з обмеженим людським контролем.


% --------------------------------------------------
\subsection{Метрики}
\label{sec:metrics}
% --------------------------------------------------

Чотири кількісні метрики операціоналізують критерії успішності архітектури.
Кожна метрика має визначений метод вимірювання, базове значення та цільовий показник.

\begin{table}[ht]
\centering
\small
\begin{tabularx}{\textwidth}{@{}l X c c@{}}
\toprule
\textbf{Метрика} & \textbf{Метод вимірювання} & \textbf{Базове значення} & \textbf{Ціль} \\
\midrule
Вартість ініціалізації (токени) &
  Вхідні токени при першому API"=виклику (включаючи кеш \texttt{CLAUDE.md}); $N{=}304$ сесії &
  30\,115 (медіана) & ${\leq}10$K \\
\addlinespace
Коефіцієнт надлишковості &
  Файли, зчитані але не редаговані / загальна кількість зчитаних файлів; $N{=}180$ сесій із ${\geq}3$ зчитуваннями &
  60\% (медіана) & ${\leq}20\%$ \\
\addlinespace
Точність витягування принципів &
  Тестовий набір записів реєстру принципів; запити отримані з edit-traces, що початково ініціювали кожен принцип &
  н/д & ${\geq}80\%$ \\
\addlinespace
Латентність оновлення (push) &
  Час від релевантного коміту до оновленого дайджесту &
  н/д (без автоматизації) & ${\leq}24$\,год \\
\bottomrule
\end{tabularx}
\caption{Метрики оцінювання з виміряними базовими значеннями та цільовими показниками.
  Базові значення ініціалізації виміряні на 304 сесіях Claude Code з датасету транскриптів~\citep{ovcharov2026edittrace}.}
\label{tab:metrics}
\end{table}

П'ята метрика --- \emph{частота корекцій витягування} (retrieval-correction edit rate, Розділ~\ref{sec:retrieval-correction}) --- має якісний характер на Фазі~1.
У міру дозрівання рівня оператора частота корекцій витягування повинна зменшуватись.
Очікувана траєкторія --- спадна крива протягом перших 6--8 тижнів розгортання, що виходить на плато, коли підсистема пам'яті покриває найбільш затребувані контексти.
У термінах бойових систем: зменшення частоти корекцій витягування означає зменшення кількості непотрібних втручань оператора, спричинених застарілим контекстом місії.


% --------------------------------------------------
\subsection{Попередні результати}
\label{sec:preliminary}
% --------------------------------------------------

На момент подання доступні три базових вимірювання, отримані з датасету транскриптів~\citep{ovcharov2026edittrace} та з історії комітів проєкту.

\paragraph{Вартість ініціалізації (Експеримент~1, $N{=}304$ сесії).}
Парсинг API"=транскриптів Claude Code з пілотного датасету дає медіану вхідних токенів першого виклику $30\,115$ (середнє $29\,693$; $\sigma{=}7\,594$; P10/P90: $18\,540/41\,774$).
З цього обсягу ${\sim}17\,600$ токенів ($59\%$) витрачаються на створення кешу для \texttt{CLAUDE.md} та системного контексту --- ці токени завантажуються безумовно для кожної сесії незалежно від задачі.
Кількість зчитувань файлів (виклики інструмента Read) в середньому становить $23{,}1$ (медіана~$14$, P90: $61$), проте $72\%$ сесій мають нуль викликів Read у фазі ініціалізації: більшість сесій починаються з команд оболонки (\texttt{git status}, \texttt{ls}), а не зі зчитування файлів.
Основна ціль зменшення підсистемою пам'яті --- автоматична вартість контексту ($T$), а не кількість зчитувань файлів.

\paragraph{Зростання CLAUDE.md (Експеримент~2, 25~комітів за 85~днів).}
Файл \texttt{CLAUDE.md} проєкту зріс із $4\,099$ до $24\,148$ символів ($152$ до $474$~рядків) за 25~комітів протягом 85~днів (17~січня --- 12~квітня~2026~р., підмножина 105"=денного операційного періоду, задокументованого в~\citep{ovcharov2026edittrace}).
Лінійна регресія кількості символів від кількості днів дає нахил $216{,}7$~симв./день при $R^2 = 0{,}87$, що підтверджує приблизно лінійне зростання, заявлене в Розділі~\ref{sec:intro}.
Кількість рядків демонструє більшу дисперсію ($R^2 = 0{,}58$) через періодичне переформатування, яке змінює щільність тексту без зміни обсягу контенту.

\paragraph{Коефіцієнт надлишковості (Експеримент~3, $N{=}180$ сесій).}
Для сесій із щонайменше трьома зчитуваннями файлів у фазі ініціалізації (перші 50~кроків) медіана коефіцієнта надлишковості контексту --- файлів, зчитаних але не редагованих у подальшому --- становить $60\%$ (середнє $56{,}7\%$; $\sigma{=}26{,}3\%$).
Розподіл є правоскошеним: $66\%$ сесій витрачають більше половини зчитувань ініціалізації марно; лише $15\%$ мають надлишковість нижче $30\%$.
Зчитування вихідного коду демонструють надлишковість $78\%$, тоді як файли пам'яті (\texttt{.claude/memory/}) --- $66{,}5\%$, що свідчить про більшу цілеспрямованість структурованого витягування з пам'яті порівняно з експлоративним зчитуванням файлів.
Коефіцієнт надлишковості корелює негативно з тривалістю сесії: короткі сесії (5--10~кроків) мають надлишковість $68{,}5\%$; довші сесії ($>$50~кроків) --- $51{,}1\%$.

Ця метрика є консервативною проксі: файл може бути зчитаний для отримання контексту без подальшого редагування, і в такому випадку він вважається <<надлишковим>>, навіть коли його вміст вплинув на міркування агента.
Базове значення $60\%$ тому, ймовірно, завищує реальну надлишковість; різниця між <<зчитано-і-відредаговано>> та <<зчитано-і-використано>> залишається для подальших досліджень.

Фаза~0 (міст <<промпт--коміт>>) розгорнута з 9~травня~2026~р. з незначним впливом на продуктивність ($15$~мс на захоплення); розподільний аналіз потребує 4--6~тижнів збору даних і буде представлений в оновленій версії.
Повна оцінка підсистеми пам'яті (точність витягування, зменшення ініціалізації з увімкненою пам'яттю) буде представлена після розгортання Фази~1.3.


% --------------------------------------------------
\subsection{Загрози валідності}
\label{sec:threats}
% --------------------------------------------------

\paragraph{Один проєкт, один практик.}
Архітектура розроблена та верифікована в контексті одного проєкту (LEX~AI) та одного практика.
Це те саме обмеження, що й у супровідній роботі~\citep[\S8]{ovcharov2026edittrace}.
Пом'якшення: архітектура є проєктно"=агностичною --- трирівнева декомпозиція та дворежимне витягування не передбачають жодних припущень щодо домену, мови програмування чи кваліфікації оператора.
Твердження про узагальнюваність відкладені до Фази~2, яка вводить багатооператорну когорту.
В контексті управління БПЛА це обмеження є прийнятним: окремий оператор із закріпленим за ним комплексом --- типовий сценарій першої лінії.

\paragraph{Змішування з дозріванням кодової бази.}
Зменшення вартості ініціалізації може бути наслідком ефектів кешування, покращення курації \texttt{CLAUDE.md} або природної стабілізації зрілої кодової бази, а не підсистеми пам'яті.
Контроль: A/B"=оцінювання шляхом увімкнення та вимкнення підсистеми пам'яті на початку сесії.
Сесії з увімкненою підсистемою пам'яті порівнюються з сесіями без неї на тому самому розподілі задач протягом того самого періоду часу.

\paragraph{Якість реєстру принципів.}
На Фазі~1 реєстр принципів курується людиною.
Точність витягування тому обмежена якістю курації, а не якістю витягування.
Якщо реєстр є неповним або містить неточні формулювання принципів, система витягування добросовісно повертатиме неточні результати.
Це обмеження визнається і адресується на Фазі~2 шляхом впровадження напівавтоматичного витягування принципів з edit-traces.


% --------------------------------------------------
\subsection{Застосовність до управління БПЛА та ситуаційних центрів}
\label{sec:metrics-mapping}
% --------------------------------------------------

Метрики платформи юридичного AI зіставляються з метриками бойових систем:

\begin{table}[ht]
\centering
\caption{Зіставлення метрик верифікації з бойовими метриками.}
\label{tab:metrics-mapping}
\begin{tabular}{@{}p{5.5cm}p{5.5cm}@{}}
\toprule
\textbf{Метрика LEX AI} & \textbf{Бойовий аналог} \\
\midrule
Токени ініціалізації ($T$) & Час підготовки оператора до управління \\
Коефіцієнт надлишковості ($W$) & Когнітивне навантаження при зміні вахти \\
Retrieval-correction signal & Непотрібні корекції = потенційно запобігнуті інциденти \\
Час передачі контексту & Час передачі управління --- метрика бойової готовності \\
\bottomrule
\end{tabular}
\end{table}
